\documentclass[../../main]{subfiles}

\ifSubfilesClassLoaded{%
    \zexternaldocument*{../1.Category-Theory/exported/aux/main}%
}{}

\begin{document}


\chapter{Additive categories}

In this chapter we introduce additive categories and functors.
We also start studying complexes and define what an homotopy between them is.

The main references for this chapter are \cite{weibelIntroductionHomologicalAlgebra1994, gelfandMethodsHomologicalAlgebra2003, aluffiAlgebraChapter02009}.

The defining trait of an additive category is that its hom-sets are not merely sets; they are abelian groups.
\begin{definition}[Additive category]\label{def:additive-category}
  A \hyperref[def:category]{category} \(\category{C}\) is \textbf{additive} if it satisfies
  \begin{enumerate}[label=Ad.\arabic*), ref=Ad.\arabic*]
    \item\label{axiom:additive-hom} For every pair of objects \(X, Y \in \category{C}\) the set \(\Hom_{\category{C}}(X, Y)\) is endowed with an abelian group structure and the composition maps are bilinear with respect to the group operation.
    \item\label{axiom:additive-zero} There exists a \hyperref[def:zero-object]{zero object} \(0 \in \category{C}\).
    \item\label{axiom:additive-coproducts} \(\category{C}\) admits finite \hyperref[def:categorical-coproduct]{coproducts}.
    \item\label{axiom:additive-products} \(\category{C}\) admits finite \hyperref[def:categorical-product]{products}.
  \end{enumerate}
\end{definition}

\begin{example}~
  \begin{itemize}
    \item The \hyperref[def:category]{category} of abelian groups \(\Ab\) is \hyperref[def:additive-category]{additive}.
          More generally, for any ring \(R\) the category of \(R\)-modules \(\Mod(R)\) is \hyperref[def:additive-category]{additive}.
    \item The \hyperref[def:category]{category} of groups \(\Grp\) is not \hyperref[def:additive-category]{additive}, since given two group homomorphisms \(f, g\) its sum is not necessarilly a homomorphism.
  \end{itemize}
\end{example}

\begin{notation}\label{def:direct-sum}
  If \(X, Y \in \category{C}\) we denote by \(X \oplus Y\) their \hyperref[def:categorical-coproduct]{coproduct} and call it their \textbf{direct sum}.
\end{notation}

In an arbitrary \hyperref[def:category]{category}, \hyperref[def:categorical-product]{products} and \hyperref[def:categorical-coproduct]{coproducts} are completely independent structures.
However, the presence of an underlying abelian group structure on hom-sets forces these two notions to collapse into one another.
\begin{lemma}\label{lemma:products-coproducts-additive-category}
  Let \(\category{C}\) be a \hyperref[def:category]{category} satisfying \ref{axiom:additive-hom} and \ref{axiom:additive-zero}.

  Suppose that \(\category{C}\) verifies
  \begin{enumerate}
    \item[Ad.5)]\label{axiom:additive-morphisms} For any pair of objects \(X, Y \in \category{C}\) there exists \(Z \in \category{C}\) and morphisms
          \begin{center}
            \begin{tikzcd}
              X \arrow[dd, "\id"'] \arrow[dr, "i_{1}"] &       & Y \arrow[dd, "\id"] \arrow[dl, "i_{2}"'] \\
              & \arrow[dl, "p_{1}"] Z \arrow[dr, "p_{2}"'] &   \\
              X &       & Y
            \end{tikzcd}
          \end{center}
          such that
          \begin{align*}
            p_{1} \circ i_{1} = \id_{X} \quad p_{1} \circ i_{2} = 0 \\
            p_{2} \circ i_{2} = \id_{Y} \quad p_{2} \circ i_{1} = 0 \\
            i_{1} \circ p_{1} + i_{2} \circ p_{2} = \id_{Z}
          \end{align*}
  \end{enumerate}
  Then conditions~\ref{axiom:additive-products},~\ref{axiom:additive-coproducts} and \hyperref[axiom:additive-morphisms]{Ad.5} are equivalent and the objects \(X \times Y, X \oplus Y\) and \(Z\) are naturally \hyperref[def:categorical-isomorphism]{isomorphic}.
\end{lemma}
\begin{proof}
  We will show that conditions~\ref{axiom:additive-coproducts} and~\hyperref[axiom:additive-morphisms]{Ad.5} are equivalent and that \(X \oplus Y \cong Z\).
  A dual argument shows that~\ref{axiom:additive-products} and~\hyperref[axiom:additive-morphisms]{Ad.5} are equivalent and that \(X \times Y \cong Z\).
  \begin{itemize}
    \item[{\ref{axiom:additive-coproducts} \(\implies\) \hyperref[axiom:additive-morphisms]{Ad.5}}] Let \(X \oplus Y\) be the \hyperref[def:categorical-coproduct]{coproduct} of \(X\) and \(Y\), that comes with morphisms \(i_{1}\colon X \to X \oplus Y\) and \(i_{2}\colon Y \to X \oplus Y\).

          By definition, for all morphisms \(f\colon X \to A\), \(g\colon Y \to B\) there exist unique morphisms \(X \oplus Y \to A, X \oplus Y \to B\) such that
          \begin{center}
            \begin{tikzcd}
              X \arrow[dr, "i_{1}"] \arrow[dd, "f"'] & & Y \arrow[dl, "i_{2}"'] \arrow[dd, "g"] \\
              & X \oplus Y \arrow[dl, dotted] \arrow[dr, dotted] & \\
              A & & B
            \end{tikzcd}
          \end{center}
          In particular if \(f = \id_{X}\) and \(g = \id_{Y}\) we recover condition \hyperref[axiom:additive-morphisms]{Ad.5}.
    \item[{\hyperref[axiom:additive-morphisms]{Ad.5} \(\implies\) \ref{axiom:additive-coproducts}}] We are going to see that \(Z\) verifies the universal property of the \hyperref[def:categorical-coproduct]{coproduct}.

          Let \(f: X \to A\) be a morphism. We want to find a morphism \(\widetilde{f}\colon Z \to A\) such that \(f = \widetilde{f} \circ i_{1}\).
          We have the following
          \begin{center}
            \begin{tikzcd}
              X \arrow[dr, "i_{1}"] \arrow[dd, "\id_{X}"'] & \\
              & Z \arrow[dl, "p_{1}"] \\
              X \arrow[d, "f"'] \\
              A
            \end{tikzcd}
          \end{center}
          Therefore if \(\widetilde{f} \coloneq f \circ p_{1}\) since \(i_{1} \circ p_{1} = \id_{X}\) we have \(\widetilde{f} \circ i_{1} = f\).

          Analogously we see that any morphism \(g: Y \to B\) can can be factored through \(Z\).
          Therefore since \(Z\) verifies the universal property of the \hyperref[def:categorical-coproduct]{coproduct}, \(Z \cong X \oplus Y\). \qedhere
  \end{itemize}
\end{proof}

\begin{remark}
  Given objects \(A_{1}, \ldots, A_{n}\) and \(B_{1}, \ldots, B_{m}\) in an \hyperref[def:additive-category]{additive category} \(\category{C}\), we can represent morphisms \(f\colon A_{1} \oplus \cdots \oplus A_{n} \to B_{1} \oplus \cdots \oplus B_{m}\) as \(n \times m\) matrices
  \begin{equation*}
    \begin{pmatrix}
      f_{11} & \cdots & f_{1n} \\
      \vdots & \ddots & \vdots \\
      f_{m1} & \cdots & f_{mn}
    \end{pmatrix}
    \text{ where } f_{k\ell}\coloneq p_{k} \circ f \circ i_{\ell}\colon A_{\ell} \to B_{k}
  \end{equation*}
  Using that \(\sum_{k} i_{k} \circ p_{k} = 1\), it follows that addition and composition of matrices obey the usual rules of matrix addition and multiplication respectively.
\end{remark}

In these \hyperref[def:category]{categories}, \hyperref[def:mono-epi-morphism]{monomorphisms} and \hyperref[def:mono-epi-morphism]{epimorphisms} have a simple characterization.
\begin{lemma} \label{lemma:mono-epi-additive}
  A morphism \(f\colon X \to Y\) in an \hyperref[def:additive-category]{additive category} is a \hyperref[def:mono-epi-morphism]{monomorphism} if and only if for every morphism \(g\colon Z \to X\)
  \begin{equation*}
    f \circ g = 0 \implies g = 0
  \end{equation*}
  Dually \(f\) is an \hyperref[def:mono-epi-morphism]{epimorphism} if and only if for every morphism \(g\colon Y \to Z\)
  \begin{equation*}
    g \circ f = 0 \implies g = 0
  \end{equation*}
\end{lemma}
\begin{proof}
  We will check the \hyperref[def:mono-epi-morphism]{monomorphism} characterization.
  \begin{itemize}
    \item[\(\implies\)] Let \(f\colon X \monoarrow Y\) be a \hyperref[def:mono-epi-morphism]{monomorphism}.
          Consider a morphism \(g\colon Z \to X\) such that \(f \circ g = 0\).
          Since \(f \circ 0 = 0\) applying left-cancellativity we have that \(g = 0\).
    \item[\(\impliedby\)] Let \(g_{1}, g_{2}\colon Z \rightrightarrows X\) be two morphisms such that \(f \circ g_{1} = f \circ g_{2}\).

          Since \(\Hom(Z, Y)\) is an abelian group we have
          \begin{equation*}
            0 = f \circ g_{1} - f \circ g_{2} = f \circ (g_{1} - g_{2})
          \end{equation*}
          Therefore by hypothesis \(g_{1} - g_{2} = 0 \iff g_{1} = g_{2}\), proving that \(f\) is a \hyperref[def:mono-epi-morphism]{monomorphism}. \qedhere
  \end{itemize}
\end{proof}

Being \hyperref[def:additive-category]{additive} also works well when considering \hyperref[def:functor-category]{functor categories}.
\begin{lemma}\label{lemma:functor-category-is-additive}
  Let \(\category{I}\) be a \hyperref[def:small-category]{small category}.
  If \(\category{C}\) is \hyperref[def:additive-category]{additive} then the \hyperref[def:functor-category]{functor category} \(\Fct{\category{I}}{\category{C}}\) is \hyperref[def:additive-category]{additive}.
\end{lemma}
\begin{proof}~
  \begin{enumerate}
    \item[\ref{axiom:additive-hom}] For any two \hyperref[def:natural-transformation]{natural transformations} \(\phi, \psi\colon \category{I} \Rightarrow \category{C}\) we can define their addition componentwise: \((\phi + \psi)_{i} = \phi_{i} + \psi_{i}\) for all \(i \in \category{I}\).
          Since \(\phi_{i}, \psi_{i} \in \Hom_{\category{C}}(-, -)\) this is well defined.
    \item[\ref{axiom:additive-zero}] Consider the zero \hyperref[def:functor]{functor}
          \begin{center}
            \begin{tikzcd}
              \llap{\(0\colon\)} \category{I} \arrow[r] & \category{C} \\[-2em]
              X \arrow[d, "f"' ""{name=source, right}] & 0 \arrow[d, "0" ""{name=target, left}] \arrow[from=source, to=target, mapsto] \\
              Y  & 0
            \end{tikzcd}
          \end{center}
          This is clearly the \hyperref[def:zero-object]{zero object} in \(\Fct{\category{I}, \category{C}}\).
    \item[\ref{axiom:additive-coproducts}, \ref{axiom:additive-products}] Direct from \autoref{prop:limits-in-functor-category}. \qedhere
  \end{enumerate}
\end{proof}

Now it is time to introduce \hyperref[def:functor]{functors} into play.
In general a \hyperref[def:functor]{functor} does not preserve the abelian group structure of the hom-sets.
Therefore the correct way to move between \hyperref[def:additive-category]{additive categories} without losing information is by using additive functors.
\begin{definition}[Additive functor]\label{def:additive-functor}
  Let \(\category{C}\), \(\category{D}\) be \hyperref[def:additive-category]{additive categories}.
  A \hyperref[def:functor]{functor} \(F\colon \category{C} \to \category{D}\) is \textbf{additive} if for all \(X, Y \in \category{C}\)
  \begin{equation*}
    \Hom_{\category{C}}(X, Y) \to \Hom_{\category{D}}(F(X), F(Y))
  \end{equation*}
  is a morphism of abelian groups.
\end{definition}

\begin{example}
  For a fixed \(R\)-module \(M\) the \hyperref[def:functor]{functor} \(\Hom_{R}(M, -)\colon \Mod(R) \to \Ab\) is \hyperref[def:additive-functor]{additive}.
\end{example}

Instead of checking that the morphism \(\Hom_{\category{C}}(X, Y) \to \Hom_{\category{D}}(F(X), F(Y))\) is an homomorphism for all objects \(X\) and \(Y\) we can use the following characterization.
\begin{proposition}\label{prop:additive-functor-characterization}
  Let \(\category{C}\), \(\category{D}\) be \hyperref[def:additive-category]{additive categories}.
  A \hyperref[def:functor]{functor} \(F\colon \category{C} \to \category{D}\) is \hyperref[def:additive-functor]{additive} if and only if it commutes with finite \hyperref[def:direct-sum]{direct sums}, i.e: if \(F(0) \cong 0\) and for all \(X, Y \in \category{C}\)
  \begin{equation*}
    F(X \oplus Y) = F(X) \oplus F(Y)
  \end{equation*}
\end{proposition}
\begin{proof}~
  \begin{itemize}
    \item[\(\implies\)] Since \(F\) is a \hyperref[def:functor]{functor} \(F(\id_{0}) = \id_{F(0)}\).
          Furthermore since \(F\) is \hyperref[def:additive-functor]{additive} if \(0_{X, Y}\) is the zero morphism between \(X, Y \in \category{C}\) then \(F(0_{X, Y}) = 0_{F(X), F(Y)}\) because group homomorphisms preserve the identity.
          In particular taking \(X = Y = 0\) since \(0_{0, 0} = \id_{0}\) we have
          \begin{equation*}
            \id_{F(0)} = F(\id_{0}) = F(0_{0, 0}) = 0_{F(0), F(0)}
          \end{equation*}
          which means that \(F(0)\) is the \hyperref[def:zero-object]{zero object} in \(\category{D}\).
          
          Let \(X, Y \in \category{C}\) be two objects and \(X \oplus Y\) be its \hyperref[def:direct-sum]{direct sum}.
          Applying \(F\) to the diagram in \autoref{lemma:products-coproducts-additive-category} defining \(Z \cong X \oplus Y\) we get
          \begin{center}
            \begin{tikzcd}
              F(X) \arrow[dd, "{\id_{F(X)}}"'] \arrow[dr, "Fi_{1}"] &       & F(Y) \arrow[dd, "{\id_{F(Y)}}"] \arrow[dl, "Fi_{2}"'] \\
              & \arrow[dl, "Fp_{1}"] F(X \oplus Y) \arrow[dr, "Fp_{2}"'] &   \\
              F(X) &       & F(Y)
            \end{tikzcd}
          \end{center}
          Note that
          \begin{align*}
            Fp_{1} \circ Fi_{1} &= F(p_{1} \circ i_{1}) = F(\id_{X}) = \id_{F(X)} & Fp_{1} \circ Fi_{2} &= F(p_{1} \circ i_{2}) = F(0) = 0 \\
            Fp_{2} \circ Fi_{2} &= F(p_{2} \circ i_{2}) = F(\id_{Y}) = \id_{F(Y)} & Fp_{2} \circ Fi_{1} &= F(p_{2} \circ i_{1}) = F(0) = 0
          \end{align*}
          Furthermore since \(F\) is \hyperref[def:additive-functor]{additive}
          \begin{equation*}
            Fp \circ Fi + Fj \circ Fq = F(p \circ i + j \circ q) = F\id_{X \oplus Y} = \id_{F(X \oplus Y)}
          \end{equation*}
          Therefore by \autoref{lemma:products-coproducts-additive-category} \(F(X \oplus Y) \cong F(X) \oplus F(Y)\).
    \item[\(\impliedby\)] Applying \autoref{lemma:products-coproducts-additive-category} to \(X \oplus X\) and \(Y \oplus Y\) we get maps \(\Delta_{X}: X \to X \oplus X\) and \(\nabla_{Y}: Y \oplus Y \to Y\) such that
          \begin{align*}
            p_{1} \circ \Delta_{X} &= \id_{X} = p_{2} \circ \Delta_{X} \\
            \nabla_{Y} \circ i_{1} &= \id_{Y} = \nabla_{Y} \circ i_{2}
          \end{align*}
          Then for any pair of morphisms \(f, g: X \to Y\) in \(\category{D}\) we have
          \begin{equation*}
            f + g = \nabla_{Y} \circ (f \oplus g) \circ \Delta_{X}
          \end{equation*}
          where \(f \oplus g\) is the unique morphism such that
          \begin{equation*}
            p_{k} \circ (f \oplus g) \circ i_{j} = \begin{cases}
              f & j = k = 1 \\
              g & j = k = 2 \\
              0 & \text{otherwise}
            \end{cases}
          \end{equation*}
          Applying \(F\) to this expression we get
          \begin{equation*}
            F(f + g) = F(\nabla_{Y} \circ (f \oplus g) \circ \Delta_{X}) = F\nabla_{Y} \circ F(f \oplus g) \circ F\Delta_{X}
          \end{equation*}
          Since by hypothesis \(F\) preserves finite \hyperref[def:direct-sum]{direct sums}
          \begin{equation*}
            F\nabla_{Y} \circ F(f \oplus g) \circ F\Delta_{X} = \nabla_{F(Y)} \circ (Ff \oplus Fg) \circ \Delta_{F(X)} = Ff + Fg
          \end{equation*}
          Therefore \(F(f + g) = Ff + Fg\), which means that \(F\) is \hyperref[def:additive-functor]{additive}. \qedhere
  \end{itemize}
\end{proof}

Being \hyperref[def:additive-functor]{additive} is preserved by \hyperref[def:adjoint-functors]{adjoints}.
\begin{proposition}\label{prop:additive-functors-and-adjunctions}
  Let \(L \dashv R\colon \category{C} \to \category{D}\) be a pair of \hyperref[def:adjoint-functors]{adjoint functors} between \hyperref[def:additive-category]{additive categories}.
  Then \(L\) is \hyperref[def:additive-functor]{additive} if and only if \(R\) is \hyperref[def:additive-functor]{additive}.
\end{proposition}
\begin{proof}
  Since \hyperref[thm:adjoints-preserve-limits]{left adjoints preserve colimits and right adjoints preserve limits} we have
  \begin{align*}
    L(X \oplus Y) &= L(X) \oplus L(Y) \\
    R(X \oplus Y) &= R(X \times Y) = R(X) \times R(Y) = R(X) \oplus R(Y)
  \end{align*}

  Note that if \(R\) is \hyperref[def:additive-functor]{additive} then \(L(0) = 0\) since
  \begin{align*}
    \Hom_{\category{C}}(L(0), 0) = \Hom_{\category{D}}(0, R(0)) = \Hom_{\category{D}}(0, 0)
  \end{align*}
  and viceversa.
  
  Therefore by \autoref{prop:additive-functor-characterization} we have the result.
\end{proof}

Being \hyperref[def:additive-category]{additive}, however, is not necessarilly preserved when taking \hyperref[def:subcategory]{full subcategories}.
For example, the \hyperref[def:subcategory]{full subcategory} of \(\Ab\) consisting of all but the trivial group is clearly not \hyperref[def:additive-category]{additive}.
This motivates the following definition.
\begin{definition}[Full additive subcategory]\label{def:full-additive-subcategory}
  Let \(\category{C}\) be an \hyperref[def:additive-category]{additive category}.

  A \hyperref[def:subcategory]{full subcategory} \(\category{D}\) of \(\category{C}\) is a \textbf{full additive subcategory} if the inclusion \hyperref[def:functor]{functor} \(\iota\colon D \hookrightarrow \category{C}\) is \hyperref[def:additive-functor]{additive}.
\end{definition}

\begin{remark}
  By \autoref{prop:additive-functor-characterization} a \hyperref[def:subcategory]{full subcategory} \(\category{D}\) of \(\category{C}\) is a \hyperref[def:full-additive-subcategory]{full additive subcategory} if and only if \(0 \in \category{D}\) and \(\category{D}\) is closed under finite \hyperref[def:direct-sum]{direct sums}.
\end{remark}


\section{Complexes}

From now until the end of the chapter, let \(\category{C}\) be an \hyperref[def:additive-category]{additive category}.

First we make explicit the notion of a ``sequence'' of objects.
\begin{definition}[Differential object]\label{def:differential-object-additive}
  A \textbf{differential object} \((X^{\bullet}, \d_{X}^{\bullet})\) in \(\category{C}\) is a sequence of objects \(X^{k} \in \category{C}\) and morphisms \(\d_{X}^{k}\) for every \(k \in \mathbb{Z}\):
  \begin{center}
    \begin{tikzcd}
      \cdots \arrow[r] & X^{k-1} \arrow[r, "\d_{X}^{k-1}"] & X^{k} \arrow[r, "\d_{X}^{k}"] & X^{k+1} \arrow[r] & \cdots
    \end{tikzcd}
  \end{center}

  A morphism of differential objects \(f^{\bullet}\colon (X^{\bullet}, \d_{X}^{\bullet}) \to (Y^{\bullet}, \d_{Y}^{\bullet})\) is a collection of morphisms \(\{f^{k}\colon X^{k} \to Y^{k}\}_{k \in \mathbb{Z}}\) such that, for all \(k \in \mathbb{Z}\), the following diagram commutes
  \begin{center}
    \begin{tikzcd}
      \cdots \arrow[r] & X^{k} \arrow[r, "\d_{X}^{k}"] \arrow[d, "f^{k}"'] & X^{k+1} \arrow[d, "f^{k+1}"] \arrow[r] & \cdots \\
      \cdots \arrow[r] & Y^{k} \arrow[r, "\d_{Y}^{k}"'] & Y^{k+1} \arrow[r] & \cdots
    \end{tikzcd}
  \end{center}

  Therefore we get a \hyperref[def:category]{category} \(\category{Diff}(\category{C})\) of differential objects in \(\category{C}\).
  By abuse of notation we will often denote differential objects \((X^{\bullet}, \d_{X}^{\bullet})\) as \(X^{\bullet}\).
\end{definition}

\begin{remark}
  We can view the \hyperref[def:differential-object-additive]{category of differential objects} \(\category{Diff}(\category{C})\) as the \hyperref[def:functor-category]{functor category} \(\Fct{\mathbb{Z}}{\category{C}}\).
  Therefore by \autoref{lemma:functor-category-is-additive} \(\category{Diff}(\category{C})\) is an \hyperref[def:additive-category]{additive category}.
\end{remark}

\begin{remark}\label{rem:identification-diff}
  For all \(n \in \mathbb{Z}\) there exists a \hyperref[def:functor-types]{fully faithful functor}
  \begin{equation*}
    \iota_{n}\colon \category{C} \to \category{Diff}(\category{C})
  \end{equation*}
  which sends an object \(X \in \category{C}\) to the \hyperref[def:differential-object-additive]{differential object}
  \begin{center}
    \begin{tikzcd}
      \cdots \arrow[r] & 0 \arrow[r] & \underbrace{X}_{\text{Position } n} \arrow[r] & 0 \arrow[r] & \cdots
    \end{tikzcd}
  \end{center}
  and a morphism \(f\colon X \to Y\) in \(\category{C}\) to
  \begin{center}
    \begin{tikzcd}
      \cdots \arrow[r] & 0 \arrow[r] \arrow[d] & X \arrow[r] \arrow[d, "f"] & 0 \arrow[r] \arrow[d] & \cdots \\
      \cdots \arrow[r] & 0 \arrow[r] & Y \arrow[r] & 0 \arrow[r] & \cdots \\
    \end{tikzcd}
  \end{center}
  Therefore we can identify an object \(X \in \category{C}\) as a \hyperref[def:differential-object-additive]{differential object} using \(\iota_{0}\).
  When the context is clear, we will directly write \(X \in \category{Diff}(\category{C})\) instead of \(\iota_{0}(X) \in \category{Diff}(\category{C})\).
\end{remark}

We will focus in a particular case of \hyperref[def:differential-object-additive]{differential objects}: those where \(\d^{k+1} \circ \d^{k} = 0\).
This choice will be better justified when we introduce exact sequences in the next chapter.

\begin{definition}[Complex]\label{def:complex}
  A \textbf{complex} in \(\category{C}\) is a \hyperref[def:differential-object-additive]{differential object} \((X^{\bullet}, \d_{X}^{\bullet})\) such that for all \(k \in \mathbb{Z}\)
  \begin{equation*}
    \d_{X}^{k+1} \circ \d_{X}^{k} = 0
  \end{equation*}
  We denote by \(\Complex{\category{C}}\) the \hyperref[def:full-additive-subcategory]{full additive subcategory} of \(\category{Diff}(\category{C})\) consisting of complexes in \(\category{C}\).
\end{definition}

Before continuing, we fix some terminology about \hyperref[def:complex]{complexes}.
\begin{definition}\label{def:bounded-complex}
  A \hyperref[def:complex]{complex} \(X^{\bullet}\) is
  \begin{enumerate}
    \item \textbf{Bounded above} if \(X^{n} = 0\) of \(n \gg 0\).
    \item \textbf{Bounded below} if \(X^{n} = 0\) of \(n \ll 0\).
    \item \textbf{Bounded} if it is bounded above and below.
  \end{enumerate}
  We denote by \(\Complex[\ast]{\category{C}}\) (\(\ast = +, -, \operatorname{b}\)) the \hyperref[def:full-additive-subcategory]{full additive subcategory} of \(\Complex{\category{C}}\) consisting of bounded above (resp. bounded below, bounded) complexes.
  To emphasize that we are working with unbounded complexes we may write \(\Complex[\operatorname{ub}]{\category{C}} \coloneq \Complex{\category{C}}\).
\end{definition}

\begin{remark}\label{rem:identification-complex}
  Note that for any \(n \in \mathbb{Z}\), the \hyperref[rem:identification-diff]{inclusion functor} \(\iota_{n}\) can be viewed as a \hyperref[def:functor]{functor}
  \begin{equation*}
    \iota_{n}\colon \category{C} \to \Complex[\ast]{\category{C}}
  \end{equation*}
  Therefore we can identify \(X \in \category{C}\) with \(\iota_{0}(X) \in \Complex[\ast]{\category{C}}\).
\end{remark}

% \begin{notation}
%   For \(a \in \mathbb{Z}\) we denote \(\Complex[\geq a]{\category{C}}\) the \hyperref[def:full-additive-subcategory]{full additive subcategory} of \(\Complex{\category{C}}\) consisting of \hyperref[def:complex]{complexes} \(X^{\bullet}\) such that \(X^{j} = 0\) for all \(j < a\).
%   We define the \hyperref[def:category]{categories} \(\Complex[\leq a]{\category{C}}\) and \(\Complex[{[a, b]}]{\category{C}}\) for \(a \leq b\) analogously.
% \end{notation}

% \begin{definition}[Shift functor]\label{def:shift-functor}
%   Let \(p \in \mathbb{Z}\).
%   The \textbf{shift functor} is the \hyperref[def:functor]{functor}
%   \begin{equation*}
%     [p]: \Complex{\category{C}} \to \Complex{\category{C}}
%   \end{equation*}
%   as
%   \begin{itemize}
%     \item For all \hyperref[def:complex]{complexes} \((X^{\bullet}, \d_{X}^{\bullet}) \in \Complex{\category{C}}\)
%           \begin{align*}
%             (X[p])^{n} \coloneq X^{n+p}  \qquad \d_{X[p]}^{n} \coloneq (-1)^{p} \d_{X}^{n+p}
%           \end{align*}
%     \item For all \hyperref[def:complex]{morphisms of complexes} \(f^{\bullet}: X^{\bullet} \to Y^{\bullet}\)
%           \begin{equation*}
%             f^{\bullet}[p]: X[p] \to Y[p],\, (f[p])^{n} \coloneq f^{n+p}
%           \end{equation*}
%   \end{itemize}
% \end{definition}

% \begin{remark}
%   Note that, for all \(p \in \mathbb{Z}\)
%   \begin{align*}
%     \d_{X[p]}^{n+1} \circ \d_{X[p]}^{n} = (-1)^{p}\d_{X}^{n+p+1} \circ (-1)^{p}\d_{X}^{n+p} = (-1)^{p} \left( \d_{X}^{(n+p) + 1} \circ \d_{X}^{n+p} \right) = (-1)^{p} 0 = 0
%   \end{align*}
%   Therefore the \hyperref[def:shift-functor]{shift functor} is well defined.
% \end{remark}


% \begin{definition}[Mapping cone]\label{def:mapping-cone}
%   Let \(f^{\bullet}: X^{\bullet} \to Y^{\bullet}\) be a \hyperref[def:complex]{morphism of complexes}.
%   The \textbf{mapping cone} of \(f\), denoted \(\operatorname{Cone}(f)\) is the \hyperref[def:complex]{complex} defined by
%   \begin{align*}
%     \operatorname{Cone}(f^{\bullet})^{n} &\coloneq (X[1])^{n} \oplus Y^{n} = X^{n+1} \oplus Y^{n} \\
%     \d_{\operatorname{Cone}(f^{\bullet})}^{n} &\coloneq \begin{pmatrix} \d_{X[1]}^{n} & 0 \\ f^{n+1} & \d_{Y}^{n} \end{pmatrix} = \begin{pmatrix} -\d_{X}^{n+1} & 0 \\ f^{n+1} & \d_{Y}^{n} \end{pmatrix}
%   \end{align*}
%   Note that there a natural \hyperref[def:complex]{morphisms of complexes}
%   \begin{equation*}
%     \alpha f^{\bullet}: Y \to \operatorname{Cone}(f^{\bullet}) \qquad \beta f^{\bullet}: \operatorname{Cone}(f^{\bullet}) \to X[1]
%   \end{equation*}
% \end{definition}

% \begin{remark}
%   Note that for any \(f^{\bullet}: X^{\bullet} \to Y^{\bullet}\)
%   \begin{align*}
%     \d_{\operatorname{Cone}(f^{\bullet})}^{n+1} \circ \d_{\operatorname{Cone}(f^{\bullet})}^{n} &= \begin{pmatrix} -\d_{X}^{n+2} & 0 \\ f^{n+2} & \d_{Y}^{n+1} \end{pmatrix} \circ \begin{pmatrix} -\d_{X}^{n+1} & 0 \\ f^{n+1} & \d_{Y}^{n} \end{pmatrix} \\
%                                                                         &= \begin{pmatrix}
%                                                                           -\d_{X}^{n+2} \circ -d_{X}^{n+1} & 0 \\
%                                                                           f^{n+2} \circ -\d_{X}^{n+1} + \d_{Y}^{n+1} \circ f^{n+1} & \d_{Y}^{n+1} \circ \d_{Y}^{n}
%                                                                         \end{pmatrix} \\
%                                                                         &= 0
%   \end{align*}
%   were we have used that \(\d_{Y}^{n+1} \circ f^{n+1} = f^{n+2} \circ \d_{X}\).
  
%   Therefore the \hyperref[def:mapping-cone]{mapping cone} is well defined.
% \end{remark}


% \begin{proposition}
%   Let \(F: \category{C} \to \category{D}\) be a \hyperref[def:additive-functor]{additive functor}.
%   Then \(F\) commutes with \hyperref[def:mapping-cone]{mapping cones}: for all morphisms \(f^{\bullet}: X^{\bullet} \to Y^{\bullet}\) in \(\Complex{\category{C}}\)
%   \begin{equation*}
%     F(\operatorname{Cone}(f^{\bullet})) \cong \operatorname{Cone}(Ff^{\bullet})
%   \end{equation*}
% \end{proposition}
% \begin{proof}
%   Since \(F\) \hyperref[prop:additive-functor-characterization]{preserves finite direct sums} for all \(n \in \mathbb{Z}\)
%   \begin{equation*}
%     F(X^{n+1} \oplus Y^{n}) \cong F(X^{n+1}) \oplus F(Y^{n})
%   \end{equation*}
%   Therefore the objects of \(F(\operatorname{Cone}(f^{\bullet}))\) and \(\operatorname{Cone}(Ff^{\bullet})\) are the same.

%   Furthermore
%   \begin{align*}
%     F \d_{\operatorname{Cone}(f^{\bullet})}^{n} &= F \begin{pmatrix} -\d_{X}^{n+1} & 0 \\ f^{n+1} & \d_{Y}^{n} \end{pmatrix} \\
%                                       &= \begin{pmatrix} F(-\d_{X}^{n+1}) & F0 \\ Ff^{n+1} & F\d_{Y}^{n} \end{pmatrix} \\
%                                       &= \begin{pmatrix} -Fd_{X}^{n+1} & 0 \\ Ff^{n+1} & F\d_{Y}^{n} \end{pmatrix} \\
%                                       &= \begin{pmatrix} -d_{F(X)}^{n+1} & 0 \\ Ff^{n+1} & \d_{F(Y)}^{n} \end{pmatrix} \\
%                                       &= \d_{\operatorname{Cone}(Ff^{\bullet})}^{n}
%   \end{align*}
% \end{proof}

% \section{Double complexes}

% \begin{definition}[Double complex]\label{def:double-complex}
%   A \textbf{double complex} in \(\category{C}\) is the data of
%   \begin{equation*}
%     \{X^{n, m}, \d_{\uparrow}^{n,m}, \d_{\to}^{n,m}\}_{n,m \in \mathbb{Z}}
%   \end{equation*}
%   where \(X^{n, m} \in \category{C}\) and the \textit{differential maps}
%   \begin{equation*}
%     \d_{\uparrow}^{n,m}: X^{n, m} \to X^{n+1, m} \qquad \d_{\to}^{n,m}: X^{n,m} \to X^{n, m+1}
%   \end{equation*}
%   satisfy
%   \begin{equation*}
%     \d_{\uparrow}^{n+1, m} \circ \d_{\uparrow}^{n, m} = \d_{\to}^{n, m+1} \circ \d_{\to}^{n, m} = 0 \qquad \d_{\to}^{n+1, m} \circ \d_{\uparrow}^{n, m} = \d_{\uparrow}^{n, m+1} \circ \d_{\to}^{n, m} 
%   \end{equation*}
%   This can be represented by a commutative diagram
%   \begin{center}
%     \begin{tikzcd}
%       & \vdots & \vdots & & \\
%       \cdots \arrow[r] & X^{n+1, m} \arrow[u] \arrow[r, "{\d_{\to}^{n+1, m}}"] & X^{n+1, m+1} \arrow[u] \arrow[r] & \cdots \\
%       \cdots \arrow[r] & X^{n, m} \arrow[u, "{\d_{\uparrow}}^{n,m}"] \arrow[r, "{\d_{\to}^{n,m}}"'] & X^{n, m+1} \arrow[u, "\d_{\uparrow}^{n, m+1}"'] \arrow[r] & \cdots \\
%       & \vdots \arrow[u] & \vdots \arrow[u] &
%     \end{tikzcd}
%   \end{center}
%   We denote double complexes by \(X^{\bullet, \bullet}\)

%   A morphism of double complexes \(f^{\bullet, \bullet}: X^{\bullet, \bullet} \to Y^{\bullet, \bullet}\) consists of maps \(\{f^{n, m}: X^{n, m} \to Y^{n, m}\}_{n, m \in \mathbb{Z}}\) such that the following cube commutes for all \(n, m \in \mathbb{Z}\):
%   \begin{center}
%     \begin{tikzcd}[row sep=huge, column sep=huge]
%       & X^{n+1, m} \arrow[rr, "{\d_{X, \to}^{n+1, m}}"] \arrow[dl, "f^{n+1, m}"'] &  & X^{n+1, m+1} \arrow[dl, "f^{n+1, m+1}"] \\
%       Y^{n+1, m} \arrow[rr, crossing over, "{\d_{Y, \to}^{n+1, m}}" near start] & & Y^{n+1, m+1} \\
%       & X^{n, m} \arrow[rr, "{\d_{X, \to}^{n, m}}"' near start] \arrow[dl, "f^{n, m}"'] \arrow[uu, "{\d_{X, \uparrow}^{n, m}}"' near end] & & X^{n, m+1} \arrow[dl, "f^{n, m+1}"] \arrow[uu, "{\d_{X, \uparrow}^{n, m+1}}"'] \\
%       Y^{n, m} \arrow[rr, "{\d_{Y, \to}^{n, m}}"'] \arrow[uu, "{\d_{Y, \uparrow}^{n, m}}"] & & Y^{n, m+1} \arrow[uu, crossing over, "{\d_{Y, \uparrow}^{n, m+1}}"' near end]
%     \end{tikzcd}
%   \end{center}

%   In this way we obtain the \hyperref[def:additive-category]{additive category} \(\Complex[2]{\category{C}}\) of double complexes in \(\category{C}\).
% \end{definition}

% \begin{remark}\label{rem:functor-double-complex-to-complex}
%   There is a \hyperref[def:functor]{functor} \(F_{I}: \Complex[2]{\category{C}} \to \Complex{\Complex{C}}\) which to a \hyperref[def:double-complex]{double complex} \(X\) associates the \hyperref[def:complex]{complex} whose objects are the rows of \(X\).
%   More precisely for all \(n \in \mathbb{Z}\) consider the \hyperref[def:complex]{complex}
%   \begin{equation*}
%     X_{I}^{n} \coloneq \{X^{n, m}, \d_{\to}^{n, m}\}_{m \in \mathbb{Z}}
%   \end{equation*}
%   The collection of morphisms \(\{\d_{\to}^{n, m}\}_{m \in \mathbb{Z}}\) defines a morphism \(\d_{I}^{n}: X_{I}^{n} \to X_{I}^{n+1}\), which trivially verifies \(\d_{I}^{n+1} \circ \d_{I}^{n} = 0\).
%   Therefore \(X_{I}\) is a \hyperref[def:complex]{complex} in \(\Complex{\category{C}}\) and we have constructed the \hyperref[def:functor]{functor} \(F_{I}\).

%   Interchainging rows and columns we get a \hyperref[def:functor]{functor} \(F_{II}: \Complex[2]{\category{C}} \to \Complex{\Complex{C}}\).
% \end{remark}


% \begin{remark}\label{rem:finiteness-condition-double-complex}
%   To make things easier we will assume that all \hyperref[def:double-complex]{double complexes} \(X^{\bullet, \bullet} \in \Complex[2]{\category{C}}\)  verify the following finiteness condition:
%   \begin{equation*}
%     \forall p \in \mathbb{Z} \quad \{ (n, m) \in \mathbb{Z}^{2} \mid X^{n, m} \neq 0 \text{ and } n + m = p\} \text{ is finite.}
%   \end{equation*}
% \end{remark}


% \begin{definition}[Total complex]\label{def:tot-complex}
%   Let \(X^{\bullet, \bullet}\) be a \hyperref[def:double-complex]{double complex}.
%   The \textbf{total complex} \(\operatorname{tot}(X)^{\bullet}\) associated to \(X^{\bullet, \bullet}\) is the \hyperref[def:complex]{complex} defined by
%   \begin{equation*}
%     \operatorname{tot}(X^{\bullet, \bullet})^{n} \coloneq \bigoplus_{p+q=n} X^{p, q} \qquad \d_{\operatorname{tot}(X)}^{n} \coloneq \sum_{p+q=n} \left( \d_{\to}^{p, q} + (-1)^{p}\d_{\uparrow}^{p, q} \right)
%   \end{equation*}
% \end{definition}

% \begin{proposition}\label{prop:tot-is-additive}
%   The \hyperref[def:differential-object]{differential object} \(\{\operatorname{tot}(X^{\bullet, \bullet})^{p}, \d_{\operatorname{tot}(X^{\bullet, \bullet})}^{p}\}_{p \in \mathbb{Z}}\) is a \hyperref[def:complex]{complex} and \(\operatorname{tot}: \Complex[2]{\category{C}} \to \Complex{\category{C}}\) is an \hyperref[def:additive-functor]{additive functor}.
% \end{proposition}
% \begin{proof}
%   To see that \(\d_{\operatorname{tot}(X^{\bullet, \bullet})}^{n+1} \circ \d_{\operatorname{tot}(X^{\bullet, \bullet})}^{n} = 0\) we look at the maps from a specific source \(X^{p, q}\) with \(p+q = n\) to \(\operatorname{tot}^{n+2}\).
%   Note that in two steps the only targets possible are
%   \begin{itemize}
%     \item[\(X^{p+2, q}\)] Reached via \(\d_{\to}^{p+1, q} \circ \d_{\to}^{p, q} = 0\).
%     \item[\(X^{p, q+2}\)] Reached via \((-1)^{p}(\d_{\uparrow}^{p, q+1} \circ \d_{\uparrow}^{p, q}) = 0\).
%     \item[\(X^{p+1, q+1}\)] There are two possible paths: \((-1)^{p+1}\d_{\uparrow}^{p+1, q} \circ \d_{\to}^{p, q}\) or \(\d_{\to}^{p, q+1} \circ (-1)^{p}\d_{\uparrow}^{p, q}\).
%           Summing these we have
%           \begin{align*}
%             (-1)^{p+1}\d_{\uparrow}^{p+1, q} \circ \d_{\to}^{p, q} + \d_{\to}^{p, q+1} \circ (-1)^{p}\d_{\uparrow}^{p, q} &= (-1)^{p} \left( -\d_{\uparrow}^{p+1, q} \circ \d_{\to}^{p, q} + \d_{\to}^{p, q+1} \circ \d_{\uparrow}^{p, q} \right) \\
%                                                                                                 &= (-1)^{p} \left( -\d_{\uparrow}^{p+1, q} \circ \d_{\to}^{p, q} + \d_{\uparrow}^{p+1, q} \circ \d_{\to}^{p, q} \right) \\
%                                                                                                 &= 0
%           \end{align*}
%   \end{itemize}
%   Therefore since all possibilities lead to the map being zero, \(\d_{\operatorname{tot}(X^{\bullet, \bullet})}^{n+1} \circ \d_{\operatorname{tot}(X^{\bullet, \bullet})}^{n} = 0\).
  
%   To see that \(\operatorname{tot}\) is \hyperref[def:additive-functor]{additive} note that
%   \begin{equation*}
%     \operatorname{tot}(0^{\bullet, \bullet}) = \bigoplus_{p+q=n} 0^{p, q} = 0^{\bullet}
%   \end{equation*}
%   Furthermore for any \(X^{\bullet, \bullet}, Y^{\bullet, \bullet} \in \Complex[2]{\category{C}}\)
%   \begin{align*}
%     \operatorname{tot}(X^{\bullet, \bullet} \oplus Y^{\bullet, \bullet}) &= \bigoplus_{p+q=n} \left( X^{p, q} \oplus Y^{p, q} \right) = \left( \bigoplus_{p+q=n} X^{p, q} \right) \oplus \left( \bigoplus_{p+q=n} Y^{p, q} \right) = \operatorname{tot}(X^{\bullet, \bullet}) \oplus \operatorname{tot}(Y^{\bullet, \bullet}) \\
%     \d_{\operatorname{tot}(X^{\bullet, \bullet} \oplus Y^{\bullet, \bullet})}^{n} &= \sum_{p+q=n} \left( \d_{X^{\bullet, \bullet} \oplus Y^{\bullet, \bullet}, \to}^{p, q} + (-1)^{p}\d_{X^{\bullet, \bullet} \oplus Y^{\bullet, \bullet}, \uparrow}^{p, q} \right) \\
%                                             &= \sum_{p+q=n} \left( \left(\d_{X^{\bullet, \bullet}, \to}^{p, q} \oplus \d_{Y^{\bullet, \bullet}, \to}^{p, q} \right) + (-1)^{p} \left(\d_{X^{\bullet, \bullet}, \uparrow}^{p, q} \oplus \d_{Y^{\bullet, \bullet}, \uparrow}^{p, q} \right)  \right) \\
%                                             &= \sum_{p+q=n} \left( \left( \d_{X^{\bullet, \bullet}, \to}^{p, q} + (-1)^{p}\d_{X^{\bullet, \bullet}, \uparrow}^{p, q} \right) \oplus \left( \d_{Y^{\bullet, \bullet}, \to}^{p, q} + (-1)^{p}\d_{Y^{\bullet, \bullet}, \uparrow}^{p, q} \right) \right) \\
%                                             &=  \left( \sum_{p+q=n} \left( \d_{X^{\bullet, \bullet}, \to}^{p, q} + (-1)^{p}\d_{X^{\bullet, \bullet}, \uparrow}^{p, q} \right) \right) \oplus \left( \sum_{p+q=n} \left( \d_{Y^{\bullet, \bullet}, \to}^{p, q} + (-1)^{p}\d_{Y^{\bullet, \bullet}, \uparrow}^{p, q} \right) \right) \\
%                                             &= \d_{\operatorname{tot}(X^{\bullet, \bullet})}^{n} \oplus \d_{\operatorname{tot}(Y^{\bullet, \bullet})}^{n}
%   \end{align*}
%   Therefore by \autoref{prop:additive-functor-characterization} \(\operatorname{tot}\) is \hyperref[def:additive-functor]{additive}.
% \end{proof}

% \begin{remark}
%   Since \hyperref[def:complex]{\(\Complex{\category{C}}\) is additive} and thus \hyperref[lemma:products-coproducts-additive-category]{finite products and coproducts coincide} we can write \(\prod\) instead of \(\bigoplus\) in the \hyperref[def:tot-complex]{\(\operatorname{tot}\) definition}.
%   We can do this because we are assuming the \hyperref[rem:finiteness-condition-double-complex]{finiteness condition}.
% \end{remark}


% \begin{remark}[Bifunctors and double complexes]
%   If \(\category{C}, \category{C'}\) and \(\category{C''}\) are \hyperref[def:additive-category]{additive categories} and \(F: \category{C} \times \category{C'} \to \category{C''}\) is an \hyperref[def:additive-functor]{additive bifunctor} (i.e: it is \hyperref[def:additive-functor]{additive} in each component).
%   Then it defines an \hyperref[def:additive-functor]{additive bifunctor}
%   \begin{equation*}
%     \Complex[2]{F}: \Complex{\category{C}} \times \Complex{\category{C'}} \to \Complex[2]{\category{C''}}
%   \end{equation*}
%   In other words, if \(X \in \Complex{\category{C}}\) and \(X' \in \Complex{\category{C'}}\) are \hyperref[def:complex]{complexes} then \(\Complex[2]{F}(X, X')\) is a \hyperref[def:double-complex]{double complex}.
% \end{remark}

% \begin{example}\label{ex:tensor-product-double-complex}
%   Let \(A\) be a ring.
%   Consider the tensor product bifunctor
%   \begin{equation*}
%     - \otimes -: \Mod{A^{\op}} \times \Mod{A} \to \Ab
%   \end{equation*}
%   In what follows we will write \(\otimes\) instead of \(\Complex[2]{\otimes}\).
%   Then for any \hyperref[def:complex]{complexes} \(X^{\bullet} \in \Complex{\Mod{A^{\op}}}, Y^{\bullet} \in \Complex{\Mod{A}}\) we have
%   \begin{align*}
%     (X^{\bullet} \otimes Y^{\bullet})^{n, m} = X^{n} \otimes Y^{m} \quad \d_{\uparrow}^{n, m} = \d_{X^{\bullet}}^{n} \otimes \id_{Y^{m}} \quad \d_{\to}^{n, m} = \id_{X^{n}} \otimes \d_{Y^{\bullet}}^{m}
%   \end{align*}

%   Therefore applying \hyperref[def:tot-complex]{\(\operatorname{tot}\)} we have
%   \begin{align*}
%     (\operatorname{tot}(X, Y))^{k} &= \bigoplus_{n+m=k} X^{n} \oplus Y^{m} \\
%     \d_{\operatorname{tot}(X^{\bullet}, Y^{\bullet})}|_{X^{n} \oplus Y^{m}} &= \d_{\to}^{n, m} + (-1)^{n}\d_{\uparrow}^{n, m}
%   \end{align*}
% \end{example}


% \begin{example}[\(\Hom^{\bullet, \bullet}\) complex]\label{ex:hom-double-complex}
%   Consider the bifunctor
%   \begin{equation*}
%     \Hom_{\category{C}}: \opcategory{C} \times \category{C} \to \Ab
%   \end{equation*}
%   In what follows we will write \(\Hom_{\category{C}}^{\bullet, \bullet}\) instead of \(\Complex[2]{\Hom_{\category{C}}}\).
%   Then for any \hyperref[def:complex]{complexes} \(X^{\bullet} \in \Complex{\opcategory{C}}, Y^{\bullet} \in \Complex{\category{C}}\) we have
%   \begin{align*}
%     \Hom_{\category{C}}^{\bullet, \bullet}(X^{\bullet}, Y^{\bullet})^{n, m} = \Hom_{\category{C}}(X^{-m}, Y^{n}) \\
%     \d_{\uparrow}^{n, m} = \Hom_{\category{C}}(X^{-m}, \d_{Y^{\bullet}}^{n}) \quad \d_{\to}^{n, m} = \Hom_{\category{C}}((-1)^{m+1}\d_{X^{\bullet}}^{m+1}, Y^{n})
%   \end{align*}

%   Using that \(\Ab\) admits \hyperref[def:categorical-product]{countable products} we set
%   \begin{equation*}
%     \Hom_{\category{C}}^{\bullet}(X^{\bullet}, Y^{\bullet}) \coloneq \operatorname{tot} \Hom_{\category{C}}^{\bullet, \bullet}(X^{\bullet}, Y^{\bullet}) \in \Complex{\Ab}
%   \end{equation*}
%   Therefore we have
%   \begin{align*}
%     \Hom_{\category{C}}(X^{\bullet}, Y^{\bullet})^{n} &= \prod_{j \in \mathbb{Z}} \Hom_{\category{C}}(X^{j}, Y^{n-j}) \\
%     (\d^{n}f)^{j} &= \d_{Y}^{j+n} \circ f^{n} + (-1)^{n+1} f^{j+1} \circ \d_{X}^{j}
%   \end{align*}

%   Note that for \(X^{\bullet}, Y^{\bullet}, Z^{\bullet} \in \Complex{\category{C}}\) we have \textit{composition} map
%   \begin{equation*}
%     \Hom_{\category{C}}^{\bullet}(X^{\bullet}, Y^{\bullet}) \times \Hom_{\category{C}}(Y^{\bullet}, Z^{\bullet}) \xrightarrow{\circ} \Hom_{\category{C}}(X^{\bullet}, Z^{\bullet})
%   \end{equation*}
%   associated to the maps
%   \begin{align*}
%     \Hom_{\category{C}}(X^{\bullet}, Y^{\bullet})^{m} \times \Hom_{\category{C}}(Y^{\bullet}, Z^{\bullet})^{n} &\to \Hom_{\category{C}}(X^{\bullet}, Z^{\bullet})^{n+m} \\
%     \prod_{j \in \mathbb{Z}} \Hom_{\category{C}}(X^{j}, Y^{j+j}) \times \prod_{j \in \mathbb{Z}} \Hom_{\category{C}}(Y^{j+m}, Z^{j+m+n}) &\mapsto  \prod_{j \in \mathbb{Z}} \Hom_{\category{C}}(X^{j}, Z^{j+m+n}) 
%   \end{align*}
% \end{example}

\section{Homotopy of complexes}

As it turns out, \hyperref[def:categorical-isomorphism]{isomorphisms} between \hyperref[def:complex]{complexes} are unnecessarilly strong for our needs and we can do with a less restrictive relation between them.
\begin{definition}[Homotopy]\label{def:homotopy}~
  \begin{enumerate}
    \item A \hyperref[def:complex]{morphism of complexes} \(f^{\bullet}\colon X^{\bullet} \to Y^{\bullet}\) is \textbf{homotopic to zero} if for all \(p \in \mathbb{Z}\) there exists a morphism \(s^{p}\colon X^{p} \to Y^{p-1}\) in \(\category{C}\) such that
          \begin{equation*}
            f^{p} = s^{p+1} \circ \d_{X}^{p} + \d_{Y}^{p-1} \circ s^{p}
          \end{equation*}
          Visually
          \begin{center}
            \begin{tikzcd}
              \cdots \arrow[r] & X^{p-1} \arrow[r] \arrow[d, "f^{p-1}"'] & X^{p} \arrow[r] \arrow[d, "f^{p}"'] \arrow[dl, "s^{p}"', dotted] & X^{p+1} \arrow[d, "f^{p+1}"] \arrow[dl, "s^{p+1}"', dotted] \arrow[r] & \cdots \\
              \cdots \arrow[r] & Y^{p-1} \arrow[r] & Y^{p} \arrow[r] & Y^{p+1} \arrow[r] & \cdots
            \end{tikzcd}
          \end{center}
    \item Two morphisms \(f^{\bullet}, g^{\bullet}\colon X^{\bullet} \to Y^{\bullet}\) are \textbf{homotopic} if \(f^{\bullet}-g^{\bullet}\) is homotopic to zero.
          We denote this by \(f^{\bullet} \sim g^{\bullet}\).
    \item A \hyperref[def:complex]{complex} \(X^{\bullet} \in \Complex{\category{C}}\) is \textbf{homotopic to zero} if \(\id_{X^{\bullet}} \sim 0\).
    \item A \hyperref[def:complex]{morphism of complexes} \(f^{\bullet}\colon X^{\bullet} \to Y^{\bullet}\) is a \textbf{homotopy equivalence} if there exists \(g^{\bullet}\colon Y^{\bullet} \to X^{\bullet}\) such that \(g^{\bullet} \circ f^{\bullet} \sim \id_{X^{\bullet}}\)  and \(f^{\bullet} \circ g^{\bullet} \sim \id_{Y^{\bullet}}\).
          In this case we say that \(X^{\bullet}\) and \(Y^{\bullet}\) are \textbf{homotopy equivalent} and write \(X^{\bullet} \sim Y^{\bullet}\).
  \end{enumerate}
\end{definition}

This relation is well behaved with respect to both composition and the action of an \hyperref[def:additive-functor]{additive functor}.

\begin{lemma}\label{lemma:composition-homotopic-to-zero}
  Let \(f^{\bullet}\colon X^{\bullet} \to Y^{\bullet}\) and \(g^{\bullet}\colon Y^{\bullet} \to Z^{\bullet}\) be two \hyperref[def:complex]{morphisms of complexes} in \(\Complex{\category{C}}\).
  If either \(f^{\bullet}\) or \(g^{\bullet}\) is \hyperref[def:homotopy]{homotopic to 0} then \(g^{\bullet} \circ f^{\bullet}\) is \hyperref[def:homotopy]{homotopic to 0}.
\end{lemma}
\begin{proof}
  Assume that \(f^{\bullet}\) is \hyperref[def:homotopy]{homotopic to 0}.
  The proof consists on chasing the following diagram:
  \begin{center}
    \begin{tikzcd}
      X^{p-1} \arrow[r] \arrow[d, "f^{p-1}"'] & X^{p} \arrow[r] \arrow[d, "f^{p}"'] \arrow[dl, "s^{p}"'] & X^{p+1} \arrow[d, "f^{p+1}"] \arrow[dl, "s^{p+1}"'] \\
      Y^{p-1} \arrow[r] \arrow[d, "g^{p-1}"'] & Y^{p} \arrow[r] \arrow[d, "g^{p}"'] & X^{p+1} \arrow[d, "g^{p+1}"] \\
      Z^{p-1} \arrow[r] & Z^{p} \arrow[r] & Z^{p+1}
    \end{tikzcd}
  \end{center}

  We have that for all \(p \in \mathbb{Z}\) there exists \(s^{p}: X^{p} \to Y^{p-1}\) such that
  \begin{equation*}
    f^{p} = s^{p+1} \circ \d_{X}^{p} + \d_{Y}^{p-1} \circ s^{p}
  \end{equation*}
  Therefore
  \begin{align*}
    g^{p} \circ f^{p} &= g^{p} \circ \left( s^{p+1} \circ \d_{X}^{p} + \d_{Y}^{p-1} \circ s^{p} \right) \\
                  &= \left(g^{p} \circ s^{p+1}\right) \circ \d_{X}^{p} + \left(g^{p} \circ \d_{Y}^{p-1}\right) \circ s^{p} \\
                  &= \left(g^{p} \circ s^{p+1}\right) \circ \d_{X}^{p} + \d_{Z}^{p-1} \circ \left( g^{p-1} \circ s^{p} \right) \qedhere
  \end{align*}
\end{proof}

\begin{lemma}\label{lemma:additive-preserves-homotopy}
  Let \(F\colon \category{C} \to \category{D}\) be an \hyperref[def:additive-functor]{additive functor}.
  If \(f^{\bullet}\colon X^{\bullet} \to Y^{\bullet}\) is \hyperref[def:homotopy]{homotopic to zero} then \(Ff^{\bullet}\) is also \hyperref[def:homotopy]{homotopic to zero}.
\end{lemma}
\begin{proof}
  Since \(f^{\bullet}\) is \hyperref[def:homotopy]{homotopic to zero} for all \(p \in \mathbb{Z}\) there exists a morphism \(s^{p}\colon X^{p} \to Y^{p-1}\) such that
  \begin{equation*}
    f^{p} = s^{p+1} \circ \d_{X}^{p} + \d_{Y}^{p-1} \circ s^{p}
  \end{equation*}
  Therefore
  \begin{equation*}
    Ff^{p} = F(s^{p+1} \circ \d_{X}^{p} + \d_{Y}^{p-1} \circ s^{p}) = Fs^{p+1} \circ \d_{F(X)}^{p} + \d_{F(Y)}^{p-1} \circ Fs^{p} \qedhere
  \end{equation*}
\end{proof}


% \begin{definition}[Homotopy category]\label{def:homotopy-category}
%   Let \(\category{C}\) be a \hyperref[def:additive-category]{additive-category}.
%   For every pair of \hyperref[def:complex]{complexes} \(X^{\bullet}, Y^{\bullet} \in \Complex{\category{C}}\) set
%   \begin{equation*}
%     \operatorname{Ht}(X^{\bullet}, Y^{\bullet}) \coloneq \{f^{\bullet}: X^{\bullet} \to Y^{\bullet} \mid f \sim 0\}
%   \end{equation*}
%   The \textbf{homotopy category} \(\operatorname{K}(\category{C})\) of \(\category{C}\) is the \hyperref[def:category]{category} defined by
%   \begin{align*}
%     \operatorname{Ob}(\operatorname{K}(\category{C})) &\coloneq \operatorname{Ob}(\Complex{\category{C}}) \\
%     \Hom_{\operatorname{K}(\category{C})}(X^{\bullet}, Y^{\bullet}) &\coloneq \quot{\Hom_{\Complex{\category{C}}}(X^{\bullet}, Y^{\bullet})}{\operatorname{Ht}(X^{\bullet}, Y^{\bullet})}
%   \end{align*}

%   We define analogously the \hyperref[def:category]{categories} \(\operatorname{K}^{\ast}(\category{C})\) (\(\ast = +, -, \operatorname{b}\)).
  
%   Note that all of these \hyperref[def:category]{categories} are endowed with an \hyperref[def:equivalent-categories]{equivalence}, the \hyperref[def:shift-functor]{shift functor} \([1]: X \mapsto X[1]\).
% \end{definition}

% \begin{remark}\label{rem:functor-extension-to-homotopy}
%   Let \(F: \category{C} \to \category{C'}\) be an \hyperref[def:additive-functor]{additive functor}.
%   This defines a \hyperref[def:additive-functor]{additive functor} \(F: \Complex{\category{C}} \to \Complex{\category{C}}\).

%   If \(f: X \to Y\) is \hyperref[def:homotopy]{homotopic to \(0\)} in \(\category{C}\) then by \autoref{lemma:additive-preserves-homotopy} \(Ff\) is \hyperref[def:homotopy]{homotopic to \(0\)}.
%   Therefore \(F\) defines an \hyperref[def:additive-functor]{additive functor}
%   \begin{equation*}
%     \operatorname{K}(F): \operatorname{K}(\category{C}) \to \operatorname{K}(\category{C'})
%   \end{equation*}

%   Now let \(F: \category{C} \times \category{C'} \to \category{C''}\) be an \hyperref[def:additive-functor]{additive bifunctor}.
%   This defines an \hyperref[def:additive-functor]{additive bifunctor} \(F: \Complex{\category{C}} \times \Complex{\category{C'}} \to \Complex{\category{C''}}\) and by composing with the \(\operatorname{tot}\) \hyperref[def:functor]{functor} a bifunctor \(\Complex[+]{\category{C}} \times \Complex[+]{\category{C'}} \to \Complex[+]{\category{C''}}\).

%   If \(f: X \to Y\) \hyperref[def:homotopy]{homotopic to \(0\)} in \(\category{C}\) and \(X' \in \category{C'}\) then \(F(f, X')\) is \hyperref[def:homotopy]{homotopic to \(0\)}, and analogously interchanging \(\category{C}\) and \(\category{C'}\).
%   Therefore \(F\) defines an \hyperref[def:additive-functor]{additive bifunctor}
%   \begin{equation*}
%     \operatorname{K}(F): \operatorname{K}(\category{C}) \times \operatorname{K}(\category{C'}) \to \operatorname{K}(\category{C''})
%   \end{equation*}

%   By abuse
%   of notation when the context is clear we will write \(F\) instead of \(\operatorname{K}(F)\).
% \end{remark}

% \begin{proposition}\label{prop:hom-}
%   Let \(\category{C}\) be an \hyperref[def:additive-category]{additive category} and \(X^{\bullet}, Y^{\bullet} \in \Complex{\category{C}}\).
%   There are \hyperref[def:categorical-isomorphism]{isomorphisms}
%   \begin{align*}
%     Z^{0}(\Hom_{\category{C}}^{\bullet}(X, Y)) &\coloneq \ker \d_{\Hom_{\category{C}}^{\bullet}(X, Y)}^{0} \cong \Hom_{\Complex{category{C}}}(X, Y) \\
%     B^{0}(\Hom_{\category{C}}^{\bullet}(X, Y)) &\coloneq \operatorname{Im} \d_{\Hom_{\category{C}}^{\bullet}(X, Y)}^{-1} \cong \operatorname{Ht}(X, Y) \\
%     \cohomology{0}(\Hom_{\category{C}}^{\bullet}(X, Y)) &\coloneq \quot{\ker \d_{\Hom_{\category{C}}^{\bullet}(X, Y)}^{0}}{\operatorname{Im} \d_{\Hom_{\category{C}}^{\bullet}(X, Y)}^{-1}} \cong \Hom_{\operatorname{K}(\category{C})}(X, Y) \\
%   \end{align*}
% \end{proposition}
% \begin{proof}~
%   \begin{enumerate}
%     \item 
%   \end{enumerate}
% \end{proof}


\end{document}
